\SourceLeaf{453}%
\Needspace{6\baselineskip}
\[
\dfrac{3014{,}48021}{3014{,}5757}=\operatorname{num}(\log=-0{,}00001376)
\]
\noindent für die Basislänge nach \textsc{Peters} und für die oben angeführten Erddimensionen, die \textsc{Gauss} in der ‘Breitenbestimmung’ angegeben hat,
\par

\Needspace{6\baselineskip}
\[
\dfrac{3014{,}451}{3014{,}5757}=\operatorname{num}(\log=-0{,}00001797)
\]
\noindent für die Basislänge nach \textsc{Peters} und \textsc{Andrae} und für die oben angeführten Erddimensionen, die \textsc{Gauss} in der ‘Breitenbestimmung’ angegeben hat,
\par

\Needspace{6\baselineskip}
\[
\dfrac{3014{,}48021}{3014{,}5757}.\dfrac{57009{,}746}{57009{,}7584}=\operatorname{num}(\log=-0{,}00001386)
\]
\noindent für die Basislänge nach \textsc{Peters} und für \textsc{Walbeck}s Erddimensionen,
\par

\Needspace{6\baselineskip}
\[
\dfrac{3014{,}451}{3014{,}5757}.\dfrac{57009{,}746}{57009{,}7584}=\operatorname{num}(\log=-0{,}00001806)
\]
\noindent für die Basislänge nach \textsc{Peters} und \textsc{Andrae} und für \textsc{Walbeck}s Erddimensionen,
\par

\Needspace{6\baselineskip}
\[
\dfrac{3014{,}48021}{3014{,}5757}.\dfrac{57009{,}746}{57008{,}551}=\operatorname{num}(\log=-0{,}00000466)
\]
\noindent für die Basislänge nach \textsc{Peters} und für \textsc{Schmidt}s IV. Bestimmung der Erddimensionen,
\par

\Needspace{6\baselineskip}
\[
\dfrac{3014{,}451}{3014{,}5757}.\dfrac{57009{,}746}{57008{,}551}=\operatorname{num}(\log=-0{,}00000887)
\]
\noindent für die Basislänge nach \textsc{Peters} und \textsc{Andrae} und für \textsc{Schmidt}s IV. Bestimmung der Erddimensionen
\par

\Needspace{6\baselineskip}
\[
\dfrac{3014{,}48021}{3014{,}5757}.\dfrac{57009{,}746}{57013{,}109}=\operatorname{num}(\log=-0{,}00003938)
\]
\noindent für die Basislänge nach \textsc{Peters} und für \textsc{Bessel}s Erddimensionen,
\par

\Needspace{6\baselineskip}
\[
\dfrac{3014{,}451}{3014{,}5757}.\dfrac{57009{,}746}{57013{,}109}=\operatorname{num}(\log=-0{,}00004359)
\]
\noindent für die Basislänge nach \textsc{Peters} und \textsc{Andrae} und für \textsc{Bessel}s Erddimensionen.
\par

Die von \textsc{Schmidt} und von \textsc{Bessel} berechneten Erddimensionen setzen voraus, dass die Längenangabe von \textsc{Schumacher} für seine Braacker Basis zugrunde gelegt wird. Eine Neuberechnung der von ihnen berücksichtigten Gradmessungen mit dieser berichtigten Basislänge würde etwas andere Zahlen für die Erddimensionen ergeben. Da mehrere neue Gradmessungen voraussichtlich bald abgeschlossen sein werden, müssen diese Zahlen aber auch in kurzer Zeit durch bessere Bestimmungen ersetzt werden.

Die hier im Abdruck zusätzlich aus den Teilverzeichnissen aufgenommenen Koordinaten stimmen entweder mit denen im Gesamtverzeichnis überein oder beruhen auf weniger genauen Bestimmungen. sie können aber die nachfolgenden ‘Abrisse’ erläutern. In \textsc{Gauss}’ Nachlass befinden sich von den Teilverzeichnissen nur Nr. 1 bis 11. Ein neuer Vergleich ergab für mich die folgenden Berichtigungen:
\begin{center}\small\begin{tabular}{@{}p{.39\linewidth}rrp{.16\linewidth}@{}}
Im Gesamtverzeichnis steht:&\(-26619{,}9\)&\(-12689{,}5\)&Lauenberge\\
Im Teilverzeichnis (3) steht:&\(-26619{,}9\)&\(+12689{,}5\)&Lauenberge
\end{tabular}\end{center}
\begin{center}\begin{tabular}{@{}rrp{.55\linewidth}@{}}
\(-233491{,}171\)&\(+181317{,}782\)&Norden, Türmchen auf der hohen Kirche. Nr. 12.\\
\(-249451{,}172\)&\(+162280{,}007\)&Langeoog, F. J. \textsc{Pauls} Haus, östlicher Giebelstock. Nr. 12.
\end{tabular}\end{center}

Um die im Koordinatenverzeichnis angegebenen Punkte leichter wiederzuerkennen, kann man die auf dieser Vermessung beruhende ‘\textsc{Papen}sche Karte vom Königreich Hannover’ mit Vorteil nutzen.

Die Überschriften \(+\) südlich und \(+\) westlich habe ich hinzugefügt, damit der Rechner die Zeichen leichter richtig versteht.
\begin{flushright}\textsc{Schering.}\end{flushright}
\begin{center}\rule{.23\linewidth}{.3pt}\end{center}
